\section{绪论}

\subsection{研究背景}

随着互联网规模的持续扩大和网络应用的日益复杂，网络安全面临着前所未有的挑战。网络入侵检测系统（Intrusion Detection System, IDS）作为网络安全的核心防御机制之一，承担着识别异常流量、检测潜在攻击的重要职责。传统的入侵检测方法主要依赖基于规则的特征匹配和专家系统，这种方式需要人工维护庞大的规则库，对零日攻击和未知变种缺乏自适应能力，难以应对日益增长的网络威胁。

近年来，机器学习与深度学习技术的快速发展为入侵检测提供了新的解决思路~\cite{sharma2025machine,akuthota2025role}。通过从海量网络流量数据中自动学习行为模式，数据驱动的方法能够有效捕捉传统规则难以描述的复杂特征，显著提升检测的准确率和泛化能力。基于这一背景，本文以经典的NSL-KDD数据集为研究对象，系统比较多种机器学习模型在二分类与多分类入侵检测任务中的性能差异。

\subsection{NSL-KDD数据集}

NSL-KDD数据集是KDD Cup 1999数据集的改进版本，由Tavallaee等人提出~\cite{tavallaee2009}，针对原始KDD-99中存在的大量冗余记录和测试集样本分布不合理等问题进行了优化，是当前入侵检测研究中被广泛采用的基准数据集之一~\cite{rehman2025systematic}。该数据集由训练文件KDDTrain+.txt和测试文件KDDTest+.txt两部分构成，训练集包含125,973条网络连接记录，测试集包含22,544条记录。

每条记录由41个特征属性和1个类别标签组成。根据特征所描述的信息层次，可将其划分为四组：9个TCP连接基本特征，描述连接的基本属性如协议类型、服务类型、连接持续时间等；13个内容特征，刻画连接内部的用户行为模式；9个时间流量特征，反映最近2秒内与当前连接具有相同主机的统计信息；10个主机流量特征，基于100个连接窗口统计目标主机的流量行为。

数据集中的标签包含一种正常类型（normal）和四种大类攻击：拒绝服务攻击（Denial of Service, DoS）、探测攻击（Probe）、远程到本地攻击（Remote to Local, R2L）以及权限提升攻击（User to Root, U2R），共涵盖23种具体攻击类型。值得注意的是，测试集包含约14种在训练集中未出现的新型攻击样本，例如apache2、httptunnel、mailbomb、mscan、named、processtable、saint、sendmail、snmpgetattack、snmpguess、sqlattack、udpstorm、worm、xterm等。这种设计使得数据集能够有效评估模型对未知攻击的泛化能力，也成为多分类任务中准确率难以突破100\%的重要原因。

\subsection{攻击类别分布}

训练集与测试集在样本分布上呈现出明显的差异。二分类层面，训练集中正常样本占比约53.5\%，异常样本占比约46.5\%，分布相对均衡；而测试集中异常样本占比上升至56.9\%，正常样本下降至43.1\%，与真实网络环境中异常流量占比偏高的实际情况更为接近，具体的二分类标签分布如图1所示。

\begin{figure*}[t]
  \centering
  \includegraphics[width=\textwidth]{01_binary_label_distribution.png}
  \caption{训练集与测试集二分类标签分布}
  \label{fig:ch1-01}
\end{figure*}

四类攻击的样本数量差异极为悬殊。训练集中DoS攻击占据主导地位，样本数高达45,927条；Probe攻击次之，为11,656条；而R2L仅有995条，U2R更是仅52条，呈现出严重的长尾分布特征。测试集延续了这种不均衡格局，但各类别间的比例有所变化，大类攻击的整体分布情况如图2所示。

\begin{figure*}[t]
  \centering
  \includegraphics[width=\textwidth]{02_attack_category_distribution.png}
  \caption{训练集与测试集四类攻击大类分布}
  \label{fig:ch1-02}
\end{figure*}

为更细致地展示数据特点，图3给出了训练集与测试集中样本数量排名前15位的具体攻击类型。

\begin{figure*}[t]
  \centering
  \includegraphics[width=\textwidth]{03_attack_type_top15.png}
  \caption{训练集与测试集中样本数量Top-15攻击类型}
  \label{fig:ch1-03}
\end{figure*}

这种类别不平衡现象，特别是R2L和U2R样本的极度稀缺，是入侵检测建模中需要重点关注的挑战。

\subsection{研究目标与论文结构}

本文的研究目标是基于NSL-KDD数据集，构建并系统比较多种机器学习模型在网络入侵检测任务中的性能表现，重点关注二分类异常检测以及23种攻击类型的细粒度多分类识别能力（在结果分析中按五大攻击类别Normal、DoS、Probe、R2L、U2R进行分组讨论）。同时，深入分析类别不平衡、未知攻击泛化等实际应用中的关键问题，为入侵检测模型的工程落地提供参考依据。

本文共分为五章。第2章将详细介绍数据预处理流程，包括缺失值与异常值处理、类别型特征编码、数值型特征变换以及特征重要性筛选等关键步骤。